\documentclass[11pt]{article}

\usepackage[utf8]{inputenc}
\usepackage[T1]{fontenc}
\usepackage{textcomp}
\usepackage[a4paper,margin=1in]{geometry}
\usepackage{parskip}
\usepackage{hyperref}

% Reviewer/editor comments are set in italics; responses follow in normal type.
\newenvironment{italic}{\itshape}{\par}

\title{Response to the Editor and Reviewers\\[0.5em]
\large ITMO Regulations to Ratchet up Ambition}
\author{}
\date{}

\begin{document}

\maketitle

\noindent I thank the Editor and both Reviewers for their fast, careful and constructive assessment of the manuscript. Below I reproduce each comment in italics and explain how I complied with each of them underneath.

\bigskip

\section*{Editor}

\begin{italic}
The editors of \emph{Global Environmental Change} have secured two reviews of your manuscript. Both reviewers are positive about the paper and signal its strengths and resonance. However, both also highlight a series of minor revisions to be made. Both reviewers have provided a thorough, constructive set of comments that are self-explanatory. Please also pay particular attention to Reviewer~1's request for changes to be made to the manuscript's structure. Please find their reviews below.

% We are happy to invite you to do further work on your manuscript, attending carefully to each of the concerns raised, before re-submitting it for our further consideration. In this case, we will likely approach one or both reviewers for their further opinion.

% In your re-submitted version, please provide an accompanying letter explaining in detail the responses made to each point by the reviewers and a version of your manuscript with tracked changes included in addition to a clean version. Please also consult the guide for authors carefully to ensure your manuscript is formatted correctly (e.g.\ Harvard referencing, avoiding or minimising the use of footnotes, etc.) to avoid additional work at the resubmission stage.
\end{italic}



I thank the Editor for the opportunity to revise the manuscript. I have carefully attended to each of the Reviewers' comments as detailed below. In addition to the clean version, I provide a version with tracked changes, and I have checked the formatting against the guide for authors.


\section*{Reviewer 1}

\begin{italic}
Dear Author, thank you for the opportunity to review the Perspective article entitled ``ITMO Regulations to Ratchet up Ambition''. The topic is quite relevant and with important timing, as some countries are still submitting/revising their NDCs, and this could improve next NDC submission cycles. I have, however, identified some minor aspects that could be improved before the manuscript can be considered for publication. Below, I provide my detailed comments, divided into major and minor points.
\end{italic}



I am very grateful for this positive assessment and for the helpful comments, which I address point by point below.


\subsection*{Major comments}

\begin{italic}
As a major comment, but a fairly easy one, I would suggest that the author make some changes to the overall structure of the Perspective, merging sections or changing some sections to subsections (e.g.: sections 4, 5 and 6 being merged as one section with subsections).
\end{italic}

Done: former Sections 4, 5 and 6 have been merged into a single section, ``A coalition of the willing to jointly define NDCs'', of which they now form the three subsections.




\subsection*{Minor comments}

\begin{italic}
Page 4, line 96, footnote 3: The author uses Angola as an example of a country that has not ratcheted up its ambition. However, one could also argue that the Angolan government is trying to pledge what it considers feasible after reassessing the extent to which the country has been able to reduce emissions and approach its 2025 target. I would suggest keeping the wording as country-neutral as possible, in order to strengthen support for the author's proposal rather than emphasising perceived shortcomings in specific national contexts.
\end{italic}

Done: I removed the footnote discussing Angola. The text is now country-neutral, referring only to Climate Action Tracker's list of countries that have not ratcheted up their targets.



\begin{italic}
Page 4, line 125, ``Until now, I conflated conditional and unconditional NDC targets.'': I dare say that the author would make a stronger argument if this part, and the sequential paragraph, appeared before, as this can be a decisive argument to Global South countries.
\end{italic}

I agree that this argument can be decisive for Global South countries. I now announce it in the introduction: a new sentence in the last paragraph of the introduction states that ``The conditional NDC target of a Global South country would supersede its unconditional target in adequacy assessments once the country's conditions in terms of climate finance are met. Hence, when a country sets its unconditional target above its projected emissions and its conditional target below them, buyers would need to provide climate finance to that country before purchasing ITMOs from it, as the seller's NDC would otherwise not be ambitious enough.'' 
The detailed discussion remains in Section 2.4, where the concepts it relies on have been introduced.


\begin{italic}
Page 6, line 198, ``usewould'', typo.
\end{italic}

Fixed.



\begin{italic}
Page 7, line 208, ``2\,\textdegree C with a 75\% probability'' \& footnote 6, could use a citation.
\end{italic}

Done: the sentence now refers to a global carbon budget achieving 2\,\textdegree C with a 67\% probability, and the footnote (stating that this is equivalent to 1.9\,\textdegree C with a 50\% probability) now cites Forster et al.\ (2025).



\begin{italic}
In all subsection titles, I suggest that the author take out the dot at the end.
\end{italic}

Done: the dots at the end of subsection titles have been removed.




\begin{italic}
At this stage, I recommend minor revisions for your Perspective paper. I hope the comments above will help strengthen the manuscript.
\end{italic}



I thank the Reviewer again for the constructive comments, which have helped strengthen the manuscript.


\section*{Reviewer 2}

\begin{italic}
The study and its methodology are transparent, well-articulated, and present a commendable proposal for a way forward. I found the study to be highly interesting and informative. However, I would appreciate it if you could clarify the assumptions underlying your proposal regarding the method of defining the ambition of the Nationally Determined Contributions (NDCs) and the grouping of the coalition.
\end{italic}



I thank the Reviewer for the positive assessment. In the revision, I clarify how the ambition of NDCs is assessed and how the baseline scenario is defined; more detailed answers follow below.


\begin{italic}
While the general approach is insightful, I think certain aspects of the methodology could be more clearly defined, particularly regarding how the Business-As-Usual (BAU) scenario is delineated. A clearer definition of BAU will enhance the understanding of the proposed methodology and the implications for future analyses.
\end{italic}

The BAU is now explicitly defined in Section 2.3, and my main proposal now relies on current policies projections instead of the BAU (see my response to Main Comment 1 below).



\subsection*{Main comments}

\begin{italic}
\textbf{1. Assumptions on BAU:} It would be beneficial to outline the key assumptions related to the Business-As-Usual (BAU) scenario in greater detail. This clarification will help readers critically evaluate the validity of your approach. I propose utilizing a current policies projection instead of the BAU scenario, as BAU is often considered a no-policy baseline, which may appear unambitious (e.g., SSP 8.5) and typically does not account for existing policies. For guidelines on defining a current policies scenario, please refer to the UNEP Emissions Gap Reports (Chapter~3), as well as studies such as den Elzen et al.\ (2019) and the latest UNEP Emissions Gap Report 2025 (Olhof et al.\ 2025). These sources present methodologies for assessing current policy scenarios, which include national independent studies, global model studies, and governmental analyses. It states: ``This scenario projects global GHG emissions assuming all currently adopted and implemented policies as of a certain specific date are realized and that no additional measures are undertaken.'' I recommend including a current policies scenario or at least referencing it in your discussion. Additionally, please elaborate on the BAU scenario and specify the sources of information for current policies.
\end{italic}

Thank you for pointing out that using the BAU as a baseline instead of current policies made the proposal less stringent. In my main proposal (Sections 4.1 and 4.2), I replaced the BAU by current policies. In Section 2.3, I added footnote 3 clarifying that the BAU is understood as a projection of emissions which accounts for recent technological developments---not as an outdated baseline such as SSP~8.5. In this section, I keep BAU in the proposal since my proposal strengthens that from Michaelowa et al. (2019) by specifying that setting an NDC below the BAU is not a sufficient condition to be allowed to sell ITMOs. The remaining occurrences of the BAU appear in the description of existing proposals (Michaelowa et al., 2019; La Hoz Theuer et al., 2019), which are framed in those terms and which I cannot alter. There, I now specify that the BAU trend is frequently updated and I added the sources you suggested (den Elzen et al., 2019; UNEP Emissions Gap Report 2025) to footnote 1 on expert projections of BAU emissions.




\begin{italic}
\textbf{2. Literature Sources on Evaluating NDC Ambition and BAU (Section 2.3):} A brief description of the reference sources utilized in the study to analyze the ambition of NDCs would be beneficial. There are various methods to assess the ambition of NDCs, as discussed in the literature, such as comparing with least-cost scenarios or equity outcomes. More references that analyze this ambition would enhance your work.
\end{italic}

Done: Section 2.3 now cites studies assessing the ambition of NDCs against least-cost scenarios and equity outcomes (H\"ohne et al., 2014; Robiou du Pont et al., 2017; van den Berg et al., 2020).


\begin{italic}
\textbf{3. Current Policies and Ambition:} It is crucial that a country adopts policies that match its ambition, as clearly explained in Nascimento et al.\ (2024). Both Nascimento et al.\ and the UNEP Emissions Gap Reports (Chapter~3) show that certain countries already have NDCs that lead to surpluses from the date of submission, such as India, Russia, and Turkey. The issue of surpluses is significant, as some countries have relatively unambitious NDCs.
\end{italic}

I agree that the issue of surpluses is significant, and my proposals are designed to address it. Note that, as long as the sum of NDCs falls short of the Paris target, a country can have a surplus and still be considered unambitious: overachieving one's target does not make that target an adequate share of the global effort. In my view, carbon trading should equalize abatement costs across countries, and NDCs should only reflect burden-sharing, not abatement activity. Under my proposals, a country overachieving an unambitious NDC would thus rightly be barred from selling ITMOs.



\begin{italic}
\textbf{4. Section 2.3 -- Line 103:} While I appreciate your proposal, it's important to note that the concept of ``equal share'' is debated in climate policy discussions. I encourage you to elaborate further on what constitutes an ``equal share'' in this context. It would be valuable to reference existing literature that examines different interpretations and applications of this concept, as a more comprehensive understanding will strengthen your argument.
\end{italic}

Done: in Section 2.3, I now cite H\"ohne et al.\ (2014), Robiou du Pont et al.\ (2017), and van den Berg et al.\ (2020), referring to the studies that examine the implications of a broad set of equity interpretations for national emission allocations.




\begin{italic}
\textbf{5. Line 149:} I respectfully disagree with the assertion made here. I believe that nearly all countries, including China, must be included in the scope of discussion if we are to meet the Paris Climate Agreement goals. It's essential to consider that China's per capita emissions are significantly higher than those of the European Union and the world average. The effectiveness of our global climate strategy hinges upon the collective ambition of these major emitters. Thus, I think the assessment of NDC ambition should reflect the realities of emissions across different countries to capture a more accurate picture of progress toward climate targets.
\end{italic}


I agree that China must be part of the collective effort. To make this clear, I added ``(including China), which should decarbonize faster and provide climate finance abroad'' after ``the ratcheting up of ambition should arguably be borne by industrialized countries''.



\subsection*{Specific comments}

\begin{italic}
\textbf{P2, line 59:} This section is somewhat undefined. The Climate Action Tracker presents current policies, which I believe should be utilized here. An alternative would involve emissions projections based on equity considerations, but this approach is not suitable for this context.
\end{italic}

In this section, I describe proposals from other papers (Michaelowa et al., 2019; La Hoz Theuer et al., 2019), so I cannot substitute current policies for the BAU in their proposals. Note that La Hoz Theuer et al.\ (2019) precisely compare simple definitions of the BAU with the Climate Action Tracker's projections. In my own proposals, I now rely on current policies (see my response to Main Comment 1).




\begin{italic}
\textbf{P3, line 66:} For instance, current policies projections can be derived from a literature assessment, as demonstrated in Chapter~3 of the UNEP Emissions Gap Report (Olhoff et al., 2025) and in publications like den Elzen et al.\ (2019). However, the data is often limited at the country level for non-G20 economies, with the exception of downscaled results, such as those from the NGFS, which may not be completely reliable.
\end{italic}


Done: I added the suggested citations (den Elzen et al., 2019; UNEP Emissions Gap Report 2025) to footnote 1 on expert projections of BAU emissions.



\begin{italic}
\textbf{Page 3, line 78:} Yes, the footnote references multiple studies. You could also mention Chapter~3 of the UNEP Emissions Gap Report, which provides an assessment of studies for UNEP, using harmonization for all studies in terms of Land Use, Land-Use Change, and Forestry (LULUCF). References to the UNEP Emissions Gap Report or den Elzen et al.\ (2019) would be appropriate. Refer to: Climate Analytics, 2026. \emph{1.5\,\textdegree C National Pathway Explorer}. Available at: \url{http://1p5ndc-pathways.climateanalytics.org}.
\end{italic}

Done: the footnote now also cites the UNEP Emissions Gap Report 2025.



\begin{italic}
\textbf{Page 4, line 114, ambition gap:} This is not, for all readers, well defined. Please refer to the references, such as the UNEP emissions gap report, but also in literature it has been defined.
\end{italic}

Done: I added after ``ambition gap'': ``understood here as the gap between the Paris target and the sum of NDCs'' and a citation to the UNEP Emissions Gap Report 2025.




\begin{italic}
\textbf{Footnote 4:} This is not the best reference. I would focus on the UNEP Emissions Gap Report (Olhoff et al.\ 2025) or Nascimento et al.\ (2024), who reached similar conclusions regarding the current policies scenario. These are original works, whereas Van den Berg et al.\ primarily presented equity outcomes and an outdated BAU projection.
\end{italic}


Done: I added the UNEP Emissions Gap Report 2025 to that footnote. I kept van den Berg et al.\ (2020) alongside it, as it is the source of the specific claim regarding India's NDC relative to the BAU.



\begin{italic}
\textbf{Line 159:} Reference multiple sources, including UNEP Emissions Gap Report (2025) and literature like Rogelj et al.\ (2023).
\end{italic}

Done: I now cite the UNEP Emissions Gap Report 2025 and Rogelj et al.\ (2023) to support the statement that the global emissions implied by current NDCs do not align with the Paris temperature target.
\\

I am very grateful for the Reviewer's comments, which have helped greatly improve the manuscript. I hope that the revisions and clarifications provided above address all concerns and strengthen the paper.

\end{document}
